\documentclass[landscape,12pt,a4paper]{article}

% ============================================================
% PACKAGES
% ============================================================
\usepackage[utf8]{inputenc}
\usepackage{graphicx}
\usepackage{xcolor}
\usepackage{geometry}
\usepackage{titlesec}
\usepackage{array}
\usepackage{enumitem}
\usepackage{float}
\usepackage{tikz}
\usepackage{ragged2e}
\usepackage{amssymb}
\usepackage{hyperref}

% ============================================================
% PAGE SETUP
% ============================================================
\geometry{
    landscape,
    a4paper,
    left=1.2cm,
    right=1.2cm,
    top=0.85cm,
    bottom=0.85cm
}

% ============================================================
% COLORS
% ============================================================
\definecolor{titleBlue}{RGB}{0,48,135}
\definecolor{darkGray}{RGB}{80,80,80}
\definecolor{calloutRed}{RGB}{210,35,35}

% ============================================================
% PARAGRAPH / LINE BREAK SETTINGS
% ============================================================
\setlength{\parindent}{0pt}
\setlength{\parskip}{0pt}
\emergencystretch=2em

% ============================================================
% SECTION FORMAT
% ============================================================
\titleformat
    {\section}
    {\Large\bfseries\color{titleBlue}}
    {}
    {0pt}
    {}

% ============================================================
% HYPERLINK SETUP
% ============================================================
\hypersetup{
    colorlinks=true,
    linkcolor=titleBlue,
    urlcolor=titleBlue
}

% ============================================================
% FOOTER
% ============================================================
\newcommand{\footertext}{
    Department of Computer Science, University of Sindh
}

% ============================================================
% DOCUMENT
% ============================================================
\begin{document}


% ============================================================
% PAGE 1: TITLE PAGE
% ============================================================
\begin{titlepage}

    \centering

    \vspace*{0.35cm}

    \includegraphics[
        width=0.95\textwidth,
        height=0.80\textheight,
        keepaspectratio
    ]{Assignment Title Page.png}

    \vfill

    {\small\color{darkGray}\footertext}

\end{titlepage}


% ============================================================
% PAGE 2: HCI SCREENSHOT
% ============================================================
\newpage

\pagenumbering{arabic}
\setcounter{page}{1}

\begin{center}

    {\Large\bfseries\color{titleBlue}
    SCREENSHOT 1: HCI FOCUS}

    \vspace{0.06cm}

    \rule{0.80\textwidth}{0.5pt}

    \vspace{0.10cm}

\end{center}


% ============================================================
% HCI SCREENSHOT + DIRECT CALLOUTS
% ============================================================
\begin{center}

\begin{tikzpicture}

% ------------------------------------------------------------
% HCI IMAGE
% ------------------------------------------------------------
\node[
    inner sep=0pt,
    outer sep=0pt
] (hciimage) at (0,0)
{
    \includegraphics[
        width=0.77\textwidth
    ]{HCI.png}
};


% ============================================================
% HCI CALLOUT 1
% SEARCH BAR
% ============================================================

\node[
    circle,
    fill=calloutRed,
    draw=white,
    line width=1.2pt,
    minimum size=0.68cm,
    inner sep=0pt,
    text=white,
    font=\bfseries\large
] (hciOne) at (-7.75,3.78) {1};

\draw[
    ->,
    >=stealth,
    very thick,
    calloutRed,
    line width=1.45pt
]
(hciOne) -- (-6.20,3.42);


% ============================================================
% HCI CALLOUT 2
% 3D NAVIGATION CONTROLS
% ============================================================

\node[
    circle,
    fill=calloutRed,
    draw=white,
    line width=1.2pt,
    minimum size=0.68cm,
    inner sep=0pt,
    text=white,
    font=\bfseries\large
] (hciTwo) at (8.55,-3.58) {2};

\draw[
    ->,
    >=stealth,
    very thick,
    calloutRed,
    line width=1.45pt
]
(hciTwo) -- (7.45,-3.55);


% ============================================================
% HCI CALLOUT 3
% INFORMATION PANEL
% ============================================================

\node[
    circle,
    fill=calloutRed,
    draw=white,
    line width=1.2pt,
    minimum size=0.68cm,
    inner sep=0pt,
    text=white,
    font=\bfseries\large
] (hciThree) at (8.55,0.70) {3};

\draw[
    ->,
    >=stealth,
    very thick,
    calloutRed,
    line width=1.45pt
]
(hciThree) -- (6.55,0.70);

\end{tikzpicture}

\end{center}


\vspace{0.04cm}


% ============================================================
% HCI EXPLANATIONS
% ============================================================
\begin{enumerate}[
    label=\textbf{\arabic*.},
    leftmargin=0.85cm,
    rightmargin=0.45cm,
    nosep,
    itemsep=0.07cm
]

    \item
    \textbf{Search Bar (Top-Left):}
    Allows users to quickly search for any location worldwide.
    The autocomplete feature provides instant suggestions,
    reducing typing effort and helping users find destinations faster.

    \item
    \textbf{3D Navigation Controls (Bottom-Right):}
    Includes zoom slider, rotate, and compass tools that let users
    explore the 3D environment from any angle. These controls
    provide clear affordances for manipulating the view.

    \item
    \textbf{Information Panel (Right-Side):}
    Displays detailed information about the selected location
    including name, description, coordinates, and save options.
    This provides immediate contextual feedback to users.

\end{enumerate}


% ============================================================
% PAGE 3: CG SCREENSHOT
% ============================================================
\newpage

\begin{center}

    {\Large\bfseries\color{titleBlue}
    SCREENSHOT 2: CG FOCUS}

    \vspace{0.06cm}

    \rule{0.80\textwidth}{0.5pt}

    \vspace{0.10cm}

\end{center}


% ============================================================
% CG SCREENSHOT + DIRECT CALLOUTS
% ============================================================
\begin{center}

\begin{tikzpicture}

% ------------------------------------------------------------
% CG IMAGE
% ------------------------------------------------------------
\node[
    inner sep=0pt,
    outer sep=0pt
] (cgimage) at (0,0)
{
    \includegraphics[
        width=0.77\textwidth
    ]{CG.png}
};


% ============================================================
% CG CALLOUT 1
% 3D MESH GEOMETRY
% ============================================================

\node[
    circle,
    fill=calloutRed,
    draw=white,
    line width=1.2pt,
    minimum size=0.68cm,
    inner sep=0pt,
    text=white,
    font=\bfseries\large
] (cgOne) at (-7.75,3.78) {1};

\draw[
    ->,
    >=stealth,
    very thick,
    calloutRed,
    line width=1.45pt
]
(cgOne) -- (-3.35,1.35);


% ============================================================
% CG CALLOUT 2
% TEXTURE MAPPING
% ============================================================

\node[
    circle,
    fill=calloutRed,
    draw=white,
    line width=1.2pt,
    minimum size=0.68cm,
    inner sep=0pt,
    text=white,
    font=\bfseries\large
] (cgTwo) at (-7.75,-3.58) {2};

\draw[
    ->,
    >=stealth,
    very thick,
    calloutRed,
    line width=1.45pt
]
(cgTwo) -- (-4.20,-2.10);


% ============================================================
% CG CALLOUT 3
% LIGHTING AND SHADOWS
% ============================================================

\node[
    circle,
    fill=calloutRed,
    draw=white,
    line width=1.2pt,
    minimum size=0.68cm,
    inner sep=0pt,
    text=white,
    font=\bfseries\large
] (cgThree) at (8.55,3.35) {3};

\draw[
    ->,
    >=stealth,
    very thick,
    calloutRed,
    line width=1.45pt
]
(cgThree) -- (3.20,-0.70);

\end{tikzpicture}

\end{center}


\vspace{0.04cm}


% ============================================================
% CG EXPLANATIONS
% ============================================================
\begin{enumerate}[
    label=\textbf{\arabic*.},
    leftmargin=0.85cm,
    rightmargin=0.45cm,
    nosep,
    itemsep=0.07cm
]

    \item
    \textbf{3D Mesh Geometry:}
    The Burj Khalifa is rendered as a detailed 3D mesh structure
    with precise geometry. This creates a realistic representation
    of the world's tallest building in 3D space.

    \item
    \textbf{Texture Mapping:}
    High-resolution satellite textures are mapped onto the ground
    surface, lake, and surrounding areas. This provides realistic
    visual context of the Dubai landscape.

    \item
    \textbf{Lighting \& Shadows:}
    Dynamic lighting simulates the sun's position, creating realistic
    shadows on the building and surrounding area. This enhances
    depth perception and makes the 3D scene more immersive.

\end{enumerate}


% ============================================================
% PAGE 4: REFLECTION QUESTIONS
% ============================================================
\newpage

\begin{center}

    {\Large\bfseries\color{titleBlue}
    REFLECTION QUESTIONS}

    \vspace{0.06cm}

    \rule{0.80\textwidth}{0.5pt}

    \vspace{0.15cm}

\end{center}


% ============================================================
% QUESTION 1
% ============================================================
\section*{Question 1:}

\textbf{
How does the Computer Graphics quality (e.g., realistic lighting
or 3D view) help or hinder the user's ability to interact with
the software?
}

\vspace{0.12cm}

\justifying

The high-quality computer graphics in Google Earth, including
realistic lighting, detailed 3D geometry, and texture mapping,
significantly enhance the user's ability to navigate and understand
spatial relationships. The visual realism makes it easier to
recognize landmarks and orient oneself in the 3D space, which
directly supports effective interaction. However, rendering these
complex graphics requires substantial processing power, which can
lead to performance lag on older or less powerful devices, hindering
smooth interaction.


\vspace{0.30cm}


% ============================================================
% QUESTION 2
% ============================================================
\section*{Question 2:}

\textbf{
What is one simple HCI change you would make to improve the
user interface?
}

\vspace{0.12cm}

\justifying

One simple but effective HCI improvement would be to add a
customizable ``Quick Actions'' toolbar that allows users to pin
their most frequently used tools. Currently, essential tools like
measuring distance, switching to street view, and toggling 3D
buildings are buried deep within menus, requiring multiple clicks
to access. A pinned toolbar would allow users to access these tools
with a single click, significantly reducing interaction cost and
improving efficiency.


\vspace{0.30cm}


% ============================================================
% GRADING TABLE
% ============================================================
\begin{table}[H]

    \centering

    \renewcommand{\arraystretch}{1.18}

    \small

    \begin{tabular}{|
        >{\bfseries}p{3.7cm}|
        p{8.4cm}|
        >{\centering\arraybackslash}p{2cm}|
    }

        \hline

        \textbf{Component}
        &
        \textbf{Requirement}
        &
        \textbf{Points}
        \\

        \hline

        HCI Annotation
        &
        Correctly identifies 3 interaction/UI elements and their
        usability purpose.
        &
        35
        \\

        \hline

        CG Annotation
        &
        Correctly identifies 3 graphics/rendering elements
        (light, mesh, camera, etc.).
        &
        35
        \\

        \hline

        Reflection
        &
        Clear, concise answers connecting graphics rendering
        to user experience.
        &
        20
        \\

        \hline

        Clarity \& Layout
        &
        Clean screenshot layout with easily readable
        callouts/pointers.
        &
        10
        \\

        \hline

    \end{tabular}

\end{table}


% ============================================================
% FOOTER
% ============================================================
\vfill

\begin{center}
    {\small\color{darkGray}\footertext}
\end{center}


% ============================================================
% END DOCUMENT
% ============================================================
\end{document}